\documentclass[11pt,a4paper]{article}
\usepackage[utf8]{inputenc}
\usepackage[T1]{fontenc}
\usepackage[portuguese]{babel}
\usepackage{xcolor}
\usepackage{hyperref}
\usepackage{geometry}
\usepackage{tcolorbox}

\geometry{margin=2.5cm}

\definecolor{primary}{RGB}{41, 128, 185}
\definecolor{secondary}{RGB}{44, 62, 80}
\definecolor{codebg}{RGB}{240, 240, 240}

\hypersetup{colorlinks=true, linkcolor=primary, urlcolor=primary}

\title{\color{secondary}\Huge\textbf{Exercício 4.1: Readiness probe}}
\author{\color{primary}DevOps com Kubernetes -- 2026}
\date{}

\begin{document}
\maketitle

\section*{\color{primary}Instruções do Exercício}
\begin{tcolorbox}[colback=blue!5!white,colframe=blue!75!black,title=Exercício 4.1: Readiness probe]
Crie uma \texttt{ReadinessProbe} para a aplicação Ping-pong. Ela deve estar pronta quando tiver uma conexão com o banco de dados. Além disso, crie outra \texttt{ReadinessProbe} para a aplicação Log output. Ela deve estar pronta quando puder receber dados da aplicação Ping-pong. Teste se a configuração funciona aplicando todos os manifestos do cluster, \textit{exceto} o statefulset do banco de dados. O status READY dos pods deve refletir a falta do banco de dados antes dele ficar disponível.
\end{tcolorbox}

\end{document}
